\subsection*{The Secret: an offering God makes medicinal\enspace\properrefs{Sec.}}

The Secret's grammar assigns both the offering and its saving effect to God.
Dedicated sacrifices are rendered to him because he granted them to be brought
for the honor of his name; he is then petitioned to make those same gifts
remedies for the offerers. Gift and giver remain distinct, but neither the
offering nor its fruit is self-authenticating. The prayer does not say that an
offering earns mercy, specify the remedy's mode, or make sacrifice a merely
subjective devotion.

Nikolaus Gihr quotes the complete Secret in his account of the Mass's value and
efficacy. He uses its movement---honor rendered to God and remedy granted to
human beings---as theological evidence for the sacrificial act's Godward and
human fruits. Gihr quotes the prayer directly but does not comment on it
word by word. Gihr's surrounding fourfold account of praise,
thanksgiving, propitiation, and petition is his own; the Secret does not
enumerate those four categories.

The prayer is also older than its 1962 placement. Old Gelasian Book~III,
section~VI transmits it in the same three-prayer set as the Collect and
Postcommunion; the Gregorian supplement, SupG~1159--1161, later transmits that
set under a historically offset post-Pentecost Sunday heading. These witnesses
establish transmission, not authorship by Gelasius or Gregory or the date of
composition. No close ancient,
patristic, or Doctoral commentary on the Secret was located in the checked
corpus.
